% ============================================================
%  Chapter 10 — Doubly Robust Estimation and Semiparametric Efficiency
%  Compile standalone:  pdflatex chapter_10.tex
%  Compile via master:  pdflatex master.tex
% ============================================================
\documentclass[master.tex]{subfiles}

\begin{document}

% ------------------------------------------------------------
%  Title block (standalone) / chapter heading (master)
% ------------------------------------------------------------
\ifSubfilesClassLoaded{%
  \begin{center}
    {\LARGE\bfseries\color{isublue}
     Chapter 10: Doubly Robust Estimation and Semiparametric Efficiency}\\[6pt]
    {\large Theory and Methods in Causal Inference}\\[4pt]
    {\normalsize Department of Statistics, Iowa State University}
  \end{center}
  \bigskip
  \hrule height 1.2pt \relax
  \smallskip
}{%
  \chapter{Doubly Robust Estimation and Semiparametric Efficiency}
}

% ------------------------------------------------------------
%  Learning objectives
% ------------------------------------------------------------
\begin{tcolorbox}[defstyle, title={Learning Objectives}]
By the end of this lecture, students should be able to:
\begin{enumerate}
  \item Derive the AIPW estimator as a bias-corrected prediction
    estimator, and explain the role of each term.
  \item Prove the double robustness property using the law of
    iterated expectations.
  \item Describe the class of augmented IPW estimators indexed by
    arbitrary control functions $b_t(x)$, verify unbiasedness for
    any $b_t$, and identify the optimal choice that minimizes total
    variance.
  \item Interpret the augmentation step as a Hilbert-space projection
    onto the orthogonal complement of the augmentation space, and
    derive the Pythagorean variance decomposition.
  \item Carry out the projection argument in the causal inference
    setting, starting from the Horvitz--Thompson estimator.
  \item Derive the two approaches to building double robustness
    directly into the outcome model: IPW-weighted regression and the
    augmented model with clever covariate; explain why including
    $\hat\pi_i^{-1}$ acts as a debiasing correction and connect this
    to TMLE \citep{vanderLaan2006targeted}.
  \item Describe internal bias calibration as an approach to achieving
    double robustness through constrained weight optimization.
  \item State the semiparametric efficiency bound for the ATE and
    explain why the AIPW estimating function is the efficient
    influence function.
  \item Identify the product-rate condition as the key sufficient
    condition for nuisance estimation error to be asymptotically
    negligible in the AIPW estimator, explain its connection to
    Neyman orthogonality, and construct a consistent variance
    estimator and Wald confidence interval from the estimated
    influence values.
\end{enumerate}
\end{tcolorbox}


% ============================================================
\section{Why Combine Outcome Regression and Weighting?}
\label{w10:sec:motivation}
% ============================================================

Under consistency, conditional exchangeability, and positivity, the
average treatment effect (ATE)
\[
  \tau = \E\{Y(1) - Y(0)\}
\]
is identified by either of the following equivalent expressions:
\[
  \tau = \E\{\mu_1(X) - \mu_0(X)\},
\]
where $\mu_t(x) = \E(Y \mid T=t,\, X=x)$ for $t \in \{0,1\}$, or
\[
  \tau
  = \E\!\left\{
      \frac{TY}{\pi(X)} - \frac{(1-T)Y}{1-\pi(X)}
    \right\},
\]
where $\pi(x) = P(T=1 \mid X=x)$ is the propensity score.

These formulas suggest two basic estimation strategies.
\begin{itemize}
  \item \textbf{Outcome regression (prediction estimator):} estimate
    $\mu_1(x)$ and $\mu_0(x)$, then average
    $\hat\mu_1(X_i)-\hat\mu_0(X_i)$ over the sample.
  \item \textbf{Inverse probability weighting (IPW):} estimate
    $\pi(x)$ and reweight observed outcomes to reconstruct the
    potential-outcome means.
\end{itemize}

Each strategy has strengths and weaknesses.  Regression-based
estimators can be statistically efficient when the outcome model is
well specified, but may be biased under misspecification.  IPW
estimators are often conceptually appealing because they directly use
the treatment model, but they can be unstable when estimated
propensity scores are close to $0$ or $1$
\citep{rosenbaum1983central}.

This chapter develops a third strategy: combine both models to
construct an estimator that is consistent when \emph{either} model is
correct, and achieves the semiparametric efficiency bound when both
are correct.
We derive the same estimator in three complementary ways: first as a
bias-corrected prediction estimator
(Section~\ref{w10:sec:prediction}), then as the optimal member of a
class of augmented estimators (Sections~\ref{w10:sec:aipw}--\ref{w10:sec:proj-causal}),
and finally as the efficient influence function in a semiparametric
model (Section~\ref{w10:sec:eif}).

% ============================================================
\section{The Prediction Estimator and Its Bias}
\label{w10:sec:prediction}
% ============================================================

To understand the source of bias it is useful to work first with
generic \emph{working regression functions} $m_t(x)$, not necessarily
equal to the true conditional means.  The \emph{prediction estimator}
(also called the regression estimator) based on $m_t$ is
\[
  \hat\tau_{\mathrm{pred}}(m)
  = \frac{1}{n}\sum_{i=1}^n
    \bigl\{m_1(X_i) - m_0(X_i)\bigr\}.
\]
Define the corresponding population prediction estimand
$\tau_{\mathrm{pred}}(m) = \E\{m_1(X)-m_0(X)\}$.
Under ignorability, the target parameter is
$\tau = \E\{\mu_1(X)-\mu_0(X)\}$,
where $\mu_t(x)=\E\{Y(t)\mid X=x\}=\E(Y\mid T=t,X=x)$.  Hence
\begin{align}
  \tau_{\mathrm{pred}}(m) - \tau
  = \E\bigl[(m_1(X)-\mu_1(X)) - (m_0(X)-\mu_0(X))\bigr].
  \label{w10:eq:pred-bias}
\end{align}
This shows that the prediction estimator is unbiased only when the
working regressions equal the true outcome regressions.  In practice
we substitute fitted models $m_t = \hat\mu_t$, and the same
decomposition explains the first-order bias induced by outcome-model
misspecification.

If the propensity score model is available, we can estimate the two
bias terms in \eqref{w10:eq:pred-bias} from the observed data.
Define residuals $e_i(t) = Y_i(t) - \hat\mu_t(X_i)$.  Since
$Y_i(1)$ is observed only for treated units and $Y_i(0)$ only for
control units,
\[
  \widehat{\mathrm{Bias}}(\hat\tau_{\mathrm{pred}})
  = -\frac{1}{n}\sum_{i=1}^n \frac{T_i}{\pi(X_i)}\,e_i(1)
    + \frac{1}{n}\sum_{i=1}^n \frac{1-T_i}{1-\pi(X_i)}\,e_i(0).
\]
The bias-corrected prediction estimator is therefore
\begin{align}
  \hat\tau_{\mathrm{AIPW}}
  = \hat\tau_{\mathrm{pred}}
    - \widehat{\mathrm{Bias}}(\hat\tau_{\mathrm{pred}})
  = \hat\mu_{1,\mathrm{dr}} - \hat\mu_{0,\mathrm{dr}},
  \label{w10:eq:aipw-biascorr}
\end{align}
where, using estimated nuisance functions $\hat\mu_t$ and
$\hat\pi$,
\begin{align}
  \hat\mu_{1,\mathrm{dr}}
  &= \frac{1}{n}\sum_{i=1}^n \frac{T_i}{\hat\pi(X_i)}\,
     \bigl\{Y_i - \hat\mu_1(X_i)\bigr\}
     + \frac{1}{n}\sum_{i=1}^n \hat\mu_1(X_i),
  \label{w10:eq:mu1dr} \\[4pt]
  \hat\mu_{0,\mathrm{dr}}
  &= \frac{1}{n}\sum_{i=1}^n \frac{1-T_i}{1-\hat\pi(X_i)}\,
     \bigl\{Y_i - \hat\mu_0(X_i)\bigr\}
     + \frac{1}{n}\sum_{i=1}^n \hat\mu_0(X_i).
  \label{w10:eq:mu0dr}
\end{align}

\begin{remark}[Two roles of the augmentation terms]
\label{w10:rem:two-interp}
The terms $T_i\{Y_i-\hat\mu_1(X_i)\}/\hat\pi(X_i)$ and
$(1-T_i)\{Y_i-\hat\mu_0(X_i)\}/(1-\hat\pi(X_i))$ have two
complementary roles.  From the bias-correction perspective, they are
IPW-weighted residuals that estimate the prediction error of the
outcome regression.  From an estimating-equation perspective
(Section~\ref{w10:sec:aipw}), they are exactly the terms that make
the estimating function orthogonal to nuisance perturbations.  The
estimated bias term has zero expectation when the outcome model is
correct; the bias-corrected estimator is unbiased when the propensity
model is correct.  Hence the name \emph{doubly robust}.
\end{remark}

% ============================================================
\section{The Augmented IPW Estimator}
\label{w10:sec:aipw}
% ============================================================

Collecting the terms in \eqref{w10:eq:aipw-biascorr} gives the
estimating-equation form.  The AIPW estimator solves
$\mathbb{P}_n\{\phi(O;\tau,\hat\eta)\} = 0$ where
\begin{align}
  \phi(O;\tau,\eta)
  &= \left[
       \frac{T}{\pi(X)}\{Y-\mu_1(X)\}+\mu_1(X)
     \right]
     - \left[
         \frac{1-T}{1-\pi(X)}\{Y-\mu_0(X)\}+\mu_0(X)
       \right]
     - \tau,
  \label{w10:eq:dr-score}
\end{align}
and $O=(X,T,Y)$ denotes one observation.  Solving explicitly,
\begin{align}
  \hat\tau_{\mathrm{AIPW}}
  = \frac{1}{n}\sum_{i=1}^n
    \Bigg[
      \hat\mu_1(X_i) - \hat\mu_0(X_i)
      + \frac{T_i}{\hat\pi(X_i)}\{Y_i - \hat\mu_1(X_i)\}
      - \frac{1-T_i}{1-\hat\pi(X_i)}\{Y_i - \hat\mu_0(X_i)\}
    \Bigg].
  \label{w10:eq:aipw}
\end{align}

\begin{defn}{Augmented Inverse Probability Weighted Estimator}
\label{w10:defn:aipw}
The estimator in \eqref{w10:eq:aipw} is the \emph{augmented
inverse probability weighted (AIPW) estimator} of the ATE.  It
equals the prediction estimator plus a weighted-residual bias
correction.
\end{defn}

The AIPW estimator enjoys two important properties: double robustness
and optimality within the augmented estimator class.  We establish
double robustness here; the optimality property is developed in
Section~\ref{w10:sec:class}.

\begin{thm}{Double Robustness}
\label{w10:thm:dr}
\label{w10:thm:dr}
Let $\phi(O;\tau,\eta)$ be the estimating function in
\eqref{w10:eq:dr-score}.  Suppose consistency, conditional
exchangeability, and positivity hold, so that
$\tau = \E\{\mu_1^*(X)-\mu_0^*(X)\}$ where
$\mu_t^*(x) = \E(Y\mid T=t,X=x)$.  Then
$\E\{\phi(O;\tau,\eta)\} = 0$ if either:
\begin{enumerate}
  \item $\mu_t(x) = \mu_t^*(x)$ for $t=0,1$, regardless of
    whether $\pi(x)$ is correctly specified; or
  \item $\pi(x) = P(T=1\mid X=x)$, regardless of whether $\mu_0$
    and $\mu_1$ are correctly specified.
\end{enumerate}
\end{thm}

\begin{proof}
\textbf{Case~1: outcome regression correct.}
If $\mu_t(X) = \E(Y\mid T=t,X)$, then conditional on $X$,
\[
  \E\!\left[\frac{T}{\pi(X)}\{Y-\mu_1(X)\}\,\middle|\,X\right]
  = \frac{P(T=1\mid X)}{\pi(X)}\,\E\{Y-\mu_1(X)\mid T=1,X\} = 0,
\]
and similarly the untreated augmentation term vanishes.  Therefore
$\E\{\phi\} = \E\{\mu_1(X)-\mu_0(X)\} - \tau = 0$.

\medskip\noindent
\textbf{Case~2: propensity score correct.}
If $\pi(X) = P(T=1\mid X)$, then conditional on $X$,
$\E[TY/\pi(X)\mid X] = \E(Y\mid T=1,X)$ and
$\E[(1-T)Y/(1-\pi(X))\mid X] = \E(Y\mid T=0,X)$, so
$\E\{\phi\mid X\} = \E(Y\mid T=1,X)-\E(Y\mid T=0,X)-\tau$.
Taking expectations over $X$ gives $\E\{\phi\} = 0$.
\end{proof}

\begin{remark}[Scope of double robustness]
\label{w10:rem:dr-scope}
Double robustness guarantees consistency if \emph{either} model is
correct; it does \emph{not} protect against misspecification of
both models simultaneously.
\end{remark}

The importance of the AIPW estimator goes beyond bias correction.  In
Section~\ref{w10:sec:class} we show that it is the variance-optimal
member of a natural class of augmented estimators indexed by arbitrary
control functions $b_t(x)$.  Later, in Section~\ref{w10:sec:eif}, we
will see that the same estimating function $\phi$ in
\eqref{w10:eq:dr-score} is exactly the efficient influence function
for the ATE under the nonparametric model.  Thus the same object
arises from three perspectives: bias correction, optimal augmentation,
and semiparametric efficiency.

% ============================================================
\section{A Class of Augmented Estimators}
\label{w10:sec:class}
% ============================================================

The AIPW estimator has a second important property beyond double
robustness: it is the optimal estimator within a broad class of
augmented IPW estimators.  We now define this class and establish
the optimality result.

For any square-integrable functions $b_1(x)$ and $b_0(x)$, define
\begin{align}
  \hat\mu_{1,b}
  &= \frac{1}{n}\sum_{i=1}^n \frac{T_i}{\pi(X_i)}\,Y_i
     - \frac{1}{n}\sum_{i=1}^n
       \left\{\frac{T_i}{\pi(X_i)}-1\right\} b_1(X_i),
  \label{w10:eq:mu1b} \\[4pt]
  \hat\mu_{0,b}
  &= \frac{1}{n}\sum_{i=1}^n \frac{1-T_i}{1-\pi(X_i)}\,Y_i
     - \frac{1}{n}\sum_{i=1}^n
       \left\{\frac{1-T_i}{1-\pi(X_i)}-1\right\} b_0(X_i),
  \label{w10:eq:mu0b}
\end{align}
and let $\hat\tau_b = \hat\mu_{1,b} - \hat\mu_{0,b}$.

\begin{remark}[Connection to the AIPW form]
\label{w10:rem:equiv}
The estimator $\hat\mu_{1,b}$ in \eqref{w10:eq:mu1b} and
$\hat\mu_{1,\mathrm{dr}}$ in \eqref{w10:eq:mu1dr} are identical
when $b_1(x) = \mu_1(x)$:
$\frac{T}{\pi}Y - (\frac{T}{\pi}-1)\mu_1
 = \frac{T}{\pi}\{Y-\mu_1\}+\mu_1$.
So the class \eqref{w10:eq:mu1b}--\eqref{w10:eq:mu0b} contains the
standard AIPW estimator as the special case $b_t = \mu_t$.
\end{remark}

The following theorem shows that every estimator in this class is
unbiased for $\tau$ under a correctly specified propensity score,
regardless of the choice of $b_t$, and provides an explicit formula
for its total variance.

\begin{thm}{Unbiasedness and Variance of the AIPW Class}
\label{w10:thm:cond-unbiased}
Under strong ignorability, SUTVA, positivity, and a correctly
specified propensity score $\pi(x) = P(T=1\mid X=x)$, for any
square-integrable $b_0$, $b_1$:
\begin{enumerate}
  \item \textbf{Unbiasedness:} $\E(\hat\tau_b) = \tau$.
  \item \textbf{Total variance:}
    \begin{align}
      \mathrm{Var}(\hat\tau_b)
      &= \frac{1}{n}\,\E\!\left[
           \left(\frac{1}{\pi(X)}-1\right)
           \{Y(1)-b_1(X)\}^2
           + 2\{Y(1)-b_1(X)\}\{Y(0)-b_0(X)\}
         \right. \nonumber \\
      &\hspace{5em}\left.
           + \left(\frac{1}{1-\pi(X)}-1\right)
           \{Y(0)-b_0(X)\}^2
         \right]
      + \frac{1}{n}\,\mathrm{Var}\{Y(1)-Y(0)\}.
      \label{w10:eq:var-b}
    \end{align}
\end{enumerate}
\end{thm}

\begin{proof}
Let $\xi = \frac{T}{\pi(X)}\{Y-b_1(X)\} + b_1(X)
           - \frac{1-T}{1-\pi(X)}\{Y-b_0(X)\} - b_0(X)$,
so that $\hat\tau_b = n^{-1}\sum_{i=1}^n \xi_i$ and, for i.i.d.\
observations, $\mathrm{Var}(\hat\tau_b) = n^{-1}\mathrm{Var}(\xi)$.
By SUTVA, $Y = TY(1)+(1-T)Y(0)$, so
\[
  \xi = \frac{T}{\pi(X)}\{Y(1)-b_1(X)\} + b_1(X)
        - \frac{1-T}{1-\pi(X)}\{Y(0)-b_0(X)\} - b_0(X).
\]

\noindent\textbf{Part~(i).}  Under strong ignorability,
$\E[T/\pi(X)\mid X] = 1$, so
$\E[(T/\pi(X)-1)b_1(X)] = \E[\E(T/\pi(X)-1\mid X)b_1(X)] = 0$.
Also $\E[T/\pi(X)\cdot Y(1)\mid X] = \E[Y(1)\mid X] = \mu_1^*(X)$.
Hence $\E[\hat\mu_{1,b}] = \mu_1$ and, by symmetry,
$\E[\hat\mu_{0,b}] = \mu_0$, so $\E(\hat\tau_b) = \tau$.

\medskip\noindent\textbf{Part~(ii).}
We apply the law of total variance, conditioning on
$(X, Y(1), Y(0))$.  Given these, the only remaining randomness
is $T \mid X \sim \mathrm{Bernoulli}(\pi(X))$.

\emph{Conditional mean.}
Since $\E[T/\pi(X)\mid X]=1$ and $\E[(1-T)/(1-\pi(X))\mid X]=1$,
\[
  \E[\xi \mid X, Y(1), Y(0)] = Y(1) - Y(0).
\]

\emph{Conditional variance.}
$\xi$ is linear in $T$:
\[
  \xi = T\!\left[
           \frac{Y(1)-b_1(X)}{\pi(X)}
           + \frac{Y(0)-b_0(X)}{1-\pi(X)}
        \right]
        - \frac{Y(0)-b_0(X)}{1-\pi(X)} + b_1(X) - b_0(X).
\]
Using $\mathrm{Var}(T\mid X) = \pi(X)(1-\pi(X))$,
\begin{align}
  \mathrm{Var}(\xi \mid X, Y(1), Y(0))
  &= \pi(X)(1\!-\!\pi(X))
     \left[\frac{Y(1)-b_1(X)}{\pi(X)}+\frac{Y(0)-b_0(X)}{1-\pi(X)}\right]^{\!2}
     \nonumber \\
  &= \left(\frac{1}{\pi(X)}-1\right)\{Y(1)-b_1(X)\}^2
     + 2\{Y(1)-b_1(X)\}\{Y(0)-b_0(X)\}
     \nonumber \\
  &\quad + \left(\frac{1}{1-\pi(X)}-1\right)\{Y(0)-b_0(X)\}^2.
  \label{w10:eq:cond-var-xi}
\end{align}

Taking expectations over $(X, Y(1), Y(0))$ and adding the between-group
term $\mathrm{Var}(\E[\xi\mid X,Y(1),Y(0)]) = \mathrm{Var}(Y(1)-Y(0))$
gives \eqref{w10:eq:var-b}.
\end{proof}

Every member of the class is unbiased for $\tau$, for any choice of
$b_t$.  The first term of \eqref{w10:eq:var-b} depends on $b_t$;
the second, $n^{-1}\mathrm{Var}\{Y(1)-Y(0)\}$, does not.

\begin{thm}{Optimal Control Functions}
\label{w10:thm:opt-b}
The functions that minimize the total variance in
\eqref{w10:eq:var-b} are
\[
  b_1^*(X) = \E\{Y(1)\mid X\}
  \qquad\text{and}\qquad
  b_0^*(X) = \E\{Y(0)\mid X\}.
\]
The resulting optimal AIPW estimator coincides with the
bias-corrected prediction estimator in \eqref{w10:eq:aipw-biascorr}.
\end{thm}

\begin{proof}
Write $\mu_t^*(x) = \E\{Y(t)\mid X=x\}$ and set
$u(X) = \mu_1^*(X)-b_1(X)$ and $v(X)=\mu_0^*(X)-b_0(X)$.
Decompose
\[
  Y(t) - b_t(X) = (\mu_t^*(X)-b_t(X)) + (Y(t)-\mu_t^*(X))
                = d_t(X) + \varepsilon_t,
\]
where $d_1=u$, $d_0=v$, and $\E[\varepsilon_t\mid X]=0$.
Under strong ignorability, $T\indep (Y(1),Y(0))\mid X$, so
$\varepsilon_t$ is also independent of $T$ given $X$.

Substituting into the $b_t$-dependent term of \eqref{w10:eq:var-b}:
\begin{align*}
  &\E\!\left[
     \left(\frac{1}{\pi}-1\right)\{Y(1)-b_1\}^2
     + 2\{Y(1)-b_1\}\{Y(0)-b_0\}
     + \left(\frac{1}{1-\pi}-1\right)\{Y(0)-b_0\}^2
   \right] \\
  &= \E\!\left[
       \left(\frac{1}{\pi}-1\right)(u+\varepsilon_1)^2
       + 2(u+\varepsilon_1)(v+\varepsilon_0)
       + \left(\frac{1}{1-\pi}-1\right)(v+\varepsilon_0)^2
     \right].
\end{align*}
Because $\E[\varepsilon_t\mid X]=0$, the law of iterated expectations
eliminates every cross term between $(u,v)$ and $(\varepsilon_1,\varepsilon_0)$:
\begin{align}
  &= \underbrace{\E\!\left[
       \left(\frac{1}{\pi(X)}-1\right)u(X)^2
       + 2\,u(X)v(X)
       + \left(\frac{1}{1-\pi(X)}-1\right)v(X)^2
     \right]}_{=:\,Q(b_0,b_1)}
     + C,
  \label{w10:eq:var-decomp}
\end{align}
where $C = \E[(1/\pi-1)\mathrm{Var}\{Y(1)\mid X\}
             +(1/(1-\pi)-1)\mathrm{Var}\{Y(0)\mid X\}]
           + 2\E[\varepsilon_1\varepsilon_0]$
does not depend on $b_0$ or $b_1$.

It remains to show $Q(b_0,b_1)\geq 0$ with equality at
$u\equiv 0$, $v\equiv 0$.  Write
\[
  \left(\frac{1}{\pi}-1\right)u^2+2uv+\left(\frac{1}{1-\pi}-1\right)v^2
  = \frac{1-\pi}{\pi}\,u^2 + 2uv + \frac{\pi}{1-\pi}\,v^2
  = \left(\sqrt{\frac{1-\pi}{\pi}}\,u+\sqrt{\frac{\pi}{1-\pi}}\,v\right)^{\!2},
\]
a perfect square, so $Q(b_0,b_1)\geq 0$.
The minimum $Q=0$ is achieved at $u\equiv 0$ and $v\equiv 0$, i.e.\
$b_t^*(X)=\mu_t^*(X)$, which gives
\[
  \mathrm{Var}(\hat\tau_b)
  \geq \frac{1}{n}\bigl[C + \mathrm{Var}\{Y(1)-Y(0)\}\bigr]
  =: \mathrm{Var}(\hat\tau_{b^*}).  \qedhere
\]
\end{proof}

\begin{remark}[Interpretation of Theorem~\ref{w10:thm:opt-b}]
\label{w10:rem:opt-b-interp}
Theorem~\ref{w10:thm:opt-b} shows that among the augmented estimators
$\{\hat\tau_b : b_t \in \mathcal{L}^2\}$, the choice
$b_t^*(X) = \E\{Y(t)\mid X\}$ minimizes asymptotic variance.  Thus
AIPW is optimal within this estimator class when both nuisance
components are correctly specified.  This is a class-specific
optimality result; the broader semiparametric efficiency
interpretation---showing that no regular asymptotically linear
estimator, inside or outside this class, can improve upon AIPW---will
be developed in Section~\ref{w10:sec:eif}.  When either model is
misspecified, the AIPW estimator retains consistency by double
robustness but no longer necessarily achieves the minimum variance
within the class.
\end{remark}

% ============================================================
\section{The Projection Interpretation}
\label{w10:sec:projection}
% ============================================================

The optimality of AIPW has an elegant interpretation in terms of
projections in a Hilbert space of estimators.
This projection view does not define a new estimator; rather, it
explains why the augmentation used in Section~\ref{w10:sec:class} is
variance reducing and why the optimal correction must be orthogonal to
the augmentation space.

\subsection{Optimal estimation via projection}
\label{w10:subsec:proj-general}

Let $\hat\theta_0$ be an unbiased estimator of $\theta$.  Define the
\emph{augmentation space}
\[
  \Lambda = \bigl\{\hat b : \E(\hat b) = 0\bigr\}
\]
as the collection of all estimators that are unbiased for zero.  For
any $\hat b \in \Lambda$ the estimator
\begin{align}
  \hat\theta_b = \hat\theta_0 - \hat b
  \label{w10:eq:theta-b}
\end{align}
remains unbiased for $\theta$.  We seek the $\hat b^* \in \Lambda$
that minimizes $\mathrm{Var}(\hat\theta_b)$.

\begin{thm}{Optimal Projection}
\label{w10:thm:projection}
The optimal correction is $\hat b^* = \Pi(\hat\theta_0\mid\Lambda)$,
the $L^2$ projection of $\hat\theta_0$ onto $\Lambda$.  It is
characterized by:
\begin{enumerate}
  \item $\hat b^* \in \Lambda$;
  \item $\mathrm{Cov}(\hat\theta_0-\hat b^*,\;\hat b) = 0$ for
    all $\hat b \in \Lambda$.
\end{enumerate}
The optimal estimator is
\[
  \hat\theta_{\mathrm{opt}}
  = \hat\theta_0 - \hat b^*
  = \Pi(\hat\theta_0 \mid \Lambda^\perp),
\]
the projection of $\hat\theta_0$ onto the orthogonal complement
$\Lambda^\perp$.  It satisfies the \emph{Pythagorean identity}
\begin{align}
  \mathrm{Var}(\hat\theta_0)
  = \mathrm{Var}(\hat\theta_{\mathrm{opt}})
    + \mathrm{Var}(\hat b^*),
  \label{w10:eq:pythagoras}
\end{align}
so in particular
$\mathrm{Var}(\hat\theta_{\mathrm{opt}}) \leq
 \mathrm{Var}(\hat\theta_0)$.
\end{thm}

The Pythagorean identity \eqref{w10:eq:pythagoras} shows that the
variance of the initial estimator decomposes orthogonally: the
optimal estimator retains only the component in $\Lambda^\perp$, and
the discarded component $\hat b^*$ accounts for the remainder.
\citet{tsiatis2006semiparametric} calls $\Lambda$ the
\emph{augmentation space} precisely because its elements are the
corrections that reduce variance.

\begin{figure}[h]
\centering
\begin{tikzpicture}[scale=1.1]
  \draw[->] (-0.3,0) -- (3.2,0) node[right] {$\Lambda$};
  \draw[->] (0,-0.3) -- (0,2.8) node[above] {$\Lambda^\perp$};
  \draw[->,thick,color=isublue] (0,0) -- (2.4,2.1)
    node[above right] {$\hat\theta_0$};
  \draw[->,thick,color=causalgreen] (0,0) -- (0,2.1)
    node[left] {$\hat\theta_{\mathrm{opt}}$};
  \draw[->,thick,color=backdoorred] (0,0) -- (2.4,0)
    node[below] {$\hat b^*$};
  \draw[dashed,gray] (2.4,0) -- (2.4,2.1);
  \draw (2.0,0) rectangle (2.4,0.4);
\end{tikzpicture}
\caption{Geometric illustration of the projection.
$\hat\theta_0$ decomposes into $\hat\theta_{\mathrm{opt}} \in
\Lambda^\perp$ and $\hat b^* \in \Lambda$.
The Pythagorean identity gives
$\mathrm{Var}(\hat\theta_0)
 = \mathrm{Var}(\hat\theta_{\mathrm{opt}})
 + \mathrm{Var}(\hat b^*)$.}
\label{w10:fig:projection}
\end{figure}

% ============================================================
\section{Projection in the Causal Inference Setting}
\label{w10:sec:proj-causal}
% ============================================================

We now specialize the projection framework to the ATE, assuming
$\pi(X)$ is known.  To align with the arm-by-arm structure of
Section~\ref{w10:sec:class}, we decompose the Horvitz--Thompson
estimator as $\hat\tau_{\mathrm{HT}} = \hat\mu_{1,\mathrm{HT}}
- \hat\mu_{0,\mathrm{HT}}$, where
\[
  \hat\mu_{1,\mathrm{HT}} = \frac{1}{n}\sum_{i=1}^n \frac{T_i Y_i}{\pi_i},
  \qquad
  \hat\mu_{0,\mathrm{HT}} = \frac{1}{n}\sum_{i=1}^n \frac{(1-T_i)Y_i}{1-\pi_i},
\]
and $\pi_i = \pi(X_i)$.

\subsection{Arm-wise augmentation spaces}
\label{w10:subsec:aug-space}

Define separate augmentation spaces for the treated and control arms:
\begin{align}
  \Lambda_1
  &= \left\{
       n^{-1}\sum_{i=1}^n
       \!\left(\frac{T_i}{\pi_i}-1\right) b_1(X_i)
       : b_1 \in \mathcal{L}^2
     \right\},
  \label{w10:eq:aug-space-1} \\[4pt]
  \Lambda_0
  &= \left\{
       n^{-1}\sum_{i=1}^n
       \!\left(\frac{1-T_i}{1-\pi_i}-1\right) b_0(X_i)
       : b_0 \in \mathcal{L}^2
     \right\}.
  \label{w10:eq:aug-space-0}
\end{align}
Because $\E\bigl[(T_i/\pi_i-1)\mid X_i\bigr]=0$ and
$\E\bigl[(1-T_i)/(1-\pi_i)-1\mid X_i\bigr]=0$, every element
of $\Lambda_1$ (resp.\ $\Lambda_0$) has expectation zero, so
$\hat\mu_{1,\mathrm{HT}} - \hat b_1$ (resp.\
$\hat\mu_{0,\mathrm{HT}} - \hat b_0$) remains unbiased for $\mu_1$
(resp.\ $\mu_0$) for any $b_t$.  The combined augmentation space
\begin{align}
  \Lambda = \Lambda_1 \oplus \Lambda_0
  \label{w10:eq:aug-space}
\end{align}
is the direct sum, whose elements are corrections that leave
$\hat\tau_{\mathrm{HT}}$ unbiased for $\tau$.

\begin{remark}[Connection to the estimator class]
\label{w10:rem:conn-class}
The estimator $\hat\mu_{1,b}$ in \eqref{w10:eq:mu1b} is exactly
$\hat\mu_{1,\mathrm{HT}}$ augmented by an element of $\Lambda_1$.
Similarly, $\hat\mu_{0,b}$ augments $\hat\mu_{0,\mathrm{HT}}$ by an
element of $\Lambda_0$.  Subtracting gives $\hat\tau_b$, the full
AIPW class of Section~\ref{w10:sec:class}.
\end{remark}

\subsection{Computing the projections}
\label{w10:subsec:proj-compute}

\begin{thm}{Arm-wise Projections onto the Augmentation Spaces}
\label{w10:thm:proj-causal}
The optimal corrections are
\begin{align}
  \Pi(\hat\mu_{1,\mathrm{HT}} \mid \Lambda_1)
  &= n^{-1}\sum_{i=1}^n
     \!\left(\frac{T_i}{\pi_i}-1\right) b_1^*(X_i),
  &b_1^*(x) &= \E\{Y(1)\mid x\},
  \label{w10:eq:proj-causal-1} \\[4pt]
  \Pi(\hat\mu_{0,\mathrm{HT}} \mid \Lambda_0)
  &= n^{-1}\sum_{i=1}^n
     \!\left(\frac{1-T_i}{1-\pi_i}-1\right) b_0^*(X_i),
  &b_0^*(x) &= \E\{Y(0)\mid x\}.
  \label{w10:eq:proj-causal-0}
\end{align}
\end{thm}

\begin{proof}
We verify the treated-arm projection; the control arm follows by
symmetry.  By Theorem~\ref{w10:thm:projection}, we must show that
\[
  R_{1i}
  := \hat\mu_{1,\mathrm{HT}} - n^{-1}\!\sum_i
     \!\left(\frac{T_i}{\pi_i}-1\right)b_1^*(X_i)
   = \frac{1}{n}\sum_{i=1}^n \frac{T_i}{\pi_i}\{Y_i - b_1^*(X_i)\}
     + \frac{1}{n}\sum_{i=1}^n b_1^*(X_i)
\]
is uncorrelated with every element of $\Lambda_1$.  By i.i.d., it
suffices to check the unit-level covariance.  Condition on
$(X_i, Y_i(1), Y_i(0))$; the only remaining randomness is
$T_i \mid X_i \sim \mathrm{Bernoulli}(\pi_i)$.  Using
$T_i Y_i = T_i Y_i(1)$ and $b_1^*(X_i) = \E\{Y_i(1)\mid X_i\}$,
\begin{align*}
  \E\!\left[
    \frac{T_i}{\pi_i}\{Y_i(1) - b_1^*(X_i)\}
    \cdot
    \left(\frac{T_i}{\pi_i}-1\right)
  \,\middle|\, X_i, Y_i(1), Y_i(0)
  \right]
  &=
  \{Y_i(1)-b_1^*(X_i)\}
  \,\E\!\left[
    \frac{T_i}{\pi_i}\!\left(\frac{T_i}{\pi_i}-1\right)
  \,\middle|\, X_i
  \right].
\end{align*}
Since $T_i^2 = T_i$,
$\E[T_i^2/\pi_i^2 \mid X_i] = 1/\pi_i$
and
$\E[T_i/\pi_i \mid X_i] = 1$, so
$\E[T_i/\pi_i\cdot(T_i/\pi_i - 1)\mid X_i]
 = 1/\pi_i - 1$.
Taking the outer expectation and using
$\E\{Y_i(1)-b_1^*(X_i)\mid X_i\} = 0$,
the whole expression vanishes.  The constant term $b_1^*(X_i)$
contributes zero covariance with $\Lambda_1$ because
$\E[(T_i/\pi_i - 1)\mid X_i] = 0$.
\end{proof}

\subsection{Recovering the AIPW estimator}
\label{w10:subsec:proj-recover}

By Theorem~\ref{w10:thm:projection}, the optimal estimators of
$\mu_1$ and $\mu_0$ are the projections onto $\Lambda_1^\perp$ and
$\Lambda_0^\perp$, respectively:
\begin{align*}
  \hat\mu_{1,\mathrm{opt}}
  &= \hat\mu_{1,\mathrm{HT}}
     - n^{-1}\!\sum_{i=1}^n\!\left(\frac{T_i}{\pi_i}-1\right)\mu_1^*(X_i)
   = \frac{1}{n}\sum_{i=1}^n
     \left[\frac{T_i}{\pi_i}\{Y_i-\mu_1^*(X_i)\}+\mu_1^*(X_i)\right], \\
  \hat\mu_{0,\mathrm{opt}}
  &= \hat\mu_{0,\mathrm{HT}}
     - n^{-1}\!\sum_{i=1}^n\!\left(\frac{1-T_i}{1-\pi_i}-1\right)\mu_0^*(X_i)
   = \frac{1}{n}\sum_{i=1}^n
     \left[\frac{1-T_i}{1-\pi_i}\{Y_i-\mu_0^*(X_i)\}+\mu_0^*(X_i)\right].
\end{align*}
Subtracting,
\[
  \hat\tau_{\mathrm{opt}}
  = \hat\mu_{1,\mathrm{opt}} - \hat\mu_{0,\mathrm{opt}},
\]
which is precisely the AIPW estimator \eqref{w10:eq:aipw} with
$b_t^*(X) = \E\{Y(t)\mid X\}$, confirming
Theorem~\ref{w10:thm:opt-b} by a different route.  The arm-wise
decomposition makes the symmetry between the treated and control
corrections transparent and shows that the asymmetric-looking
function $b^*(x)$ in the single-space formulation is an artifact of
combining two symmetric corrections into one.

% ============================================================
\section{Doubly Robust Regression: Weighted and Augmented Approaches}
\label{w10:sec:ibc}
% ============================================================

The AIPW estimator achieves double robustness by adding an explicit
bias-correction term to the prediction estimator.  An equally
important question is how to build double robustness \emph{directly
into the outcome model fit}, so that the prediction estimator is
itself doubly robust without any separate augmentation step.  This
section presents two general approaches and connects the second to
the TMLE framework.

\subsection{The internal bias calibration conditions}
\label{w10:subsec:ibc-conditions}

Let $\hat\mu_t(x)$ denote any fitted outcome model for arm $t$, and
let $\hat\pi_i = \hat\pi(X_i)$.  The prediction estimator
$\hat\tau_{\mathrm{pred}} = n^{-1}\sum_i\{\hat\mu_1(X_i)-\hat\mu_0(X_i)\}$
requires no augmentation if the IPW-weighted residuals vanish:
\begin{align}
  \sum_{i=1}^n \frac{T_i}{\hat\pi_i}
    \bigl\{Y_i - \hat\mu_1(X_i)\bigr\} &= 0,
  \label{w10:eq:ibc-cond1} \\
  \sum_{i=1}^n \frac{1-T_i}{1-\hat\pi_i}
    \bigl\{Y_i - \hat\mu_0(X_i)\bigr\} &= 0.
  \label{w10:eq:ibc-cond2}
\end{align}
We call \eqref{w10:eq:ibc-cond1}--\eqref{w10:eq:ibc-cond2} the
\emph{internal bias calibration} (IBC) conditions
\citep{firth1998robust}, following the terminology introduced in the
context of prediction estimation in survey sampling.  The word
\emph{internal} reflects the fact that these constraints are imposed
within the outcome-model fitting step itself: once they hold, no
separate bias-correction term is needed afterward.  When both IBC
conditions hold, the augmentation terms in \eqref{w10:eq:aipw} are
zero by construction, so $\hat\tau_{\mathrm{pred}} =
\hat\tau_{\mathrm{AIPW}}$ and the prediction estimator is itself
doubly robust.  The two approaches below are two ways of fitting
$\hat\mu_t$ so that the IBC conditions are satisfied by construction.

\subsection{Weighted regression approach}
\label{w10:subsec:wls}

Suppose the outcome model for arm $t$ is parameterized as
$\mu_t(X;\theta_t)$ for a finite-dimensional parameter $\theta_t$,
with the model including a constant term.  Instead of fitting
$\theta_t$ by ordinary least squares among the units with $T_i=t$,
use \emph{IPW-weighted} least squares: minimize
\begin{align}
  \sum_{i=1}^n \frac{T_i}{\hat\pi_i}
    \bigl\{Y_i - \mu_1(X_i;\theta_1)\bigr\}^2
  \label{w10:eq:wls1}
\end{align}
for $t=1$, and analogously with $(1-T_i)/(1-\hat\pi_i)$ for $t=0$.
This is essentially IPW-weighted estimation of the regression
coefficients: each unit's contribution to the least-squares criterion
is upweighted by the inverse of its probability of being in that
treatment arm.  The normal equation of \eqref{w10:eq:wls1} with
respect to the intercept component of $\theta_1$ is
\[
  \sum_{i=1}^n \frac{T_i}{\hat\pi_i}
    \bigl\{Y_i - \mu_1(X_i;\hat\theta_1)\bigr\} = 0,
\]
which is exactly the IBC condition \eqref{w10:eq:ibc-cond1}.  Hence
the IPW-weighted fitted values automatically satisfy the IBC
condition for any model containing a constant, and the resulting
prediction estimator is doubly robust; see \citet{robins1994estimation}
and \citet{bang2005doubly} for a general treatment.
The key point is not that weighted least squares creates a
fundamentally different doubly robust estimator; rather, it constructs
fitted values for which the prediction estimator algebraically equals
the AIPW estimator.

\subsection{Augmented model approach and the clever covariate}
\label{w10:subsec:augmented}

The augmented approach achieves the same calibration by modifying
the covariate set rather than the fitting weights.  Let
$\hat\mu_1^{(0)}(X_i)$ be any initial fit of the treated outcome
model (e.g., from OLS or a flexible machine-learning method).
Augment it by including the \emph{clever covariate}
$\hat\pi_i^{-1}$ \citep{vanderLaan2006targeted,vanderLaan2011targeted}
and run OLS of $Y_i$ on $\hat\mu_1^{(0)}(X_i)$ and $\hat\pi_i^{-1}$
among treated units, fitting the model
\begin{align}
  Y_i = \alpha_1 + \beta_1\hat\mu_1^{(0)}(X_i)
        + \gamma_1\hat\pi_i^{-1} + e_i(1),
  \qquad T_i = 1.
  \label{w10:eq:aug-model}
\end{align}
Denote the OLS estimates by $\hat\alpha_1$, $\hat\beta_1$,
$\hat\gamma_1$.  The fitted outcome model evaluated at all units is
\begin{align}
  \hat\mu_1(X_i)
  = \hat\alpha_1 + \hat\beta_1\hat\mu_1^{(0)}(X_i)
    + \hat\gamma_1\hat\pi_i^{-1},
  \label{w10:eq:aug-fit}
\end{align}
and the prediction estimator of $\mu_1 = \E\{Y(1)\}$ is
\begin{align}
  \hat\mu_1^{\mathrm{pred}}
  = \frac{1}{n}\sum_{i=1}^n \hat\mu_1(X_i).
  \label{w10:eq:aug-pred}
\end{align}

The normal equation of \eqref{w10:eq:aug-model} for $\hat\gamma_1$
(the coefficient on $\hat\pi_i^{-1}$) is
\begin{align}
  \sum_{i=1}^n \frac{T_i}{\hat\pi_i}
    \bigl\{Y_i - \hat\mu_1(X_i)\bigr\} = 0,
  \label{w10:eq:score-gamma}
\end{align}
which is exactly the IBC condition \eqref{w10:eq:ibc-cond1}.
Since \eqref{w10:eq:score-gamma} equals zero, we can add it to
\eqref{w10:eq:aug-pred} without changing its value, giving the AIPW
representation:
\begin{align}
  \hat\mu_1^{\mathrm{pred}}
  = \frac{1}{n}\sum_{i=1}^n \hat\mu_1(X_i)
    + \frac{1}{n}\sum_{i=1}^n
      \frac{T_i}{\hat\pi_i}
      \bigl\{Y_i - \hat\mu_1(X_i)\bigr\}.
  \label{w10:eq:aug-aipw}
\end{align}
The IPW-weighted residual sum is zero by the IBC condition, so the
prediction estimator coincides with the AIPW estimator.

The model \eqref{w10:eq:aug-model} improves upon the simpler
fixed-offset augmentation $\hat\mu_1^{(0)}(X_i) + \hat\gamma_1\hat\pi_i^{-1}$
in two respects.  First, the free coefficient $\hat\beta_1$ allows
the scale of the initial fit to be recalibrated: if
$\hat\mu_1^{(0)}$ systematically over- or under-predicts, $\hat\beta_1
\neq 1$ absorbs this global scaling error.  Second, the intercept
$\hat\alpha_1$ absorbs any residual mean shift.  The covariate
$\hat\pi_i^{-1}$ is called \emph{clever} because it is chosen so
that the fitted regression satisfies the same score equation
\eqref{w10:eq:score-gamma} that appears in the AIPW bias correction.
When the initial outcome model is misspecified, including this
covariate lets the regression absorb the component of the
model error that would otherwise appear in the augmentation term, so
the fitted prediction estimator becomes algebraically equivalent to a
debiased estimator.  Together, $(\hat\alpha_1, \hat\beta_1,
\hat\gamma_1)$ provide additional recalibration of the initial
prediction beyond what the canonical fixed-offset form achieves,
yielding a lower residual variance and hence greater efficiency than
either the fixed-offset version or plain AIPW when the OR model is
misspecified.

\begin{remark}[Behavior under correct outcome model]
\label{w10:rem:gamma-zero}
If the initial outcome model $\hat\mu_1^{(0)}$ is correctly
specified, then $\hat\alpha_1 \overset{p}{\to} 0$,
$\hat\beta_1 \overset{p}{\to} 1$, and
$\hat\gamma_1 \overset{p}{\to} 0$: the augmented model
\eqref{w10:eq:aug-model} reduces asymptotically to
$\hat\mu_1^{(0)}(X_i)$, and $\hat\mu_1^{\mathrm{pred}}$ coincides
with the standard prediction estimator.  Including the additional
covariates therefore carries no asymptotic efficiency cost when the
OR model is correct.
\end{remark}

\begin{remark}[Connection to TMLE]
\label{w10:rem:tmle}
The augmented model approach captures the same targeting idea that
underlies \emph{targeted minimum loss-based estimation} (TMLE),
introduced by \citet{vanderLaan2006targeted} and developed
comprehensively in \citet{vanderLaan2011targeted}.  The standard TMLE
targeting step uses the fixed-offset form (coefficient on
$\hat\mu_1^{(0)}$ fixed at~1), which is the minimal fluctuation
sufficient to satisfy the IBC condition.  The model
\eqref{w10:eq:aug-model} goes further by freely estimating
$\hat\beta_1$, providing additional recalibration of the initial fit
beyond what canonical TMLE achieves in its one-dimensional targeting
step.  A thorough treatment of TMLE, including its asymptotic theory
and connections to cross-fitting, is beyond the scope of these notes;
see \citet{vanderLaan2011targeted} and \citet{rose2020targeted} for
details.
\end{remark}

% ============================================================
\section{Lab: Simulation Study}
\label{w10:sec:lab}
% ============================================================

This lab compares four estimators of the ATE: the IPW estimator
$\hat\tau_{\mathrm{HT}}$, the AIPW estimator $\hat\tau_{\mathrm{AIPW}}$,
the fixed-offset augmented-model estimator $\hat\tau_{\mathrm{aug}}$,
and the improved augmented-model estimator $\hat\tau_{\mathrm{aug}^+}$
of \eqref{w10:eq:aug-model}.  The propensity score model is correctly
specified throughout; two scenarios differ in whether the outcome
regression model is correctly specified.  Full \textsf{R} code is
provided in the supplementary file \texttt{chapter10\_lab.R}.

\subsection{Simulation setup}
\label{w10:subsec:lab-setup}

Generate $n = 1000$ i.i.d.\ observations as follows.  Draw
$X_i \sim N(0,1)$, then $T_i \mid X_i \sim
\mathrm{Bernoulli}(\pi(X_i))$ with true propensity score
$\pi(x) = \mathrm{expit}(0.5x)$.  The potential outcomes are
\begin{align*}
  Y_i(1) &= 1 + X_i + 0.5X_i^2 + \varepsilon_i(1), \\
  Y_i(0) &= \phantom{1 +{}} X_i + 0.5X_i^2 + \varepsilon_i(0),
\end{align*}
with $\varepsilon_i(t) \overset{\mathrm{i.i.d.}}{\sim} N(0,1)$
independent of $(X_i, T_i)$, giving true ATE $\tau = 1$.  The
observed outcome is $Y_i = T_i Y_i(1) + (1-T_i)Y_i(0)$.  The
quadratic term $0.5X_i^2$ is the component that a linear outcome
model omits.

The propensity score is estimated by logistic regression of $T_i$
on $X_i$ (correctly specified by construction).  Two outcome
regression scenarios are considered.
\begin{itemize}
  \item \textbf{Scenario~1 (OR correct):} fit separate linear
    regressions of $Y_i$ on $(1, X_i, X_i^2)$ within each arm.
  \item \textbf{Scenario~2 (OR misspecified):} fit separate linear
    regressions of $Y_i$ on $(1, X_i)$ only, omitting the quadratic
    term.
\end{itemize}

\subsection{Estimator construction}
\label{w10:subsec:lab-estimators}

Given $\hat\pi_i = \hat\pi(X_i)$ and initial fitted values
$\hat\mu_t(X_i) \equiv \hat\mu_t^{(0)}(X_i)$, let
$c_{1i} = \hat\pi_i^{-1}$ and $c_{0i} = (1-\hat\pi_i)^{-1}$.

\medskip\noindent
\textbf{IPW estimator.}
\[
  \hat\tau_{\mathrm{HT}}
  = \frac{1}{n}\sum_{i=1}^n
    \left\{
      \frac{T_i Y_i}{\hat\pi_i}
      - \frac{(1-T_i)Y_i}{1-\hat\pi_i}
    \right\}.
\]

\medskip\noindent
\textbf{AIPW estimator.}
\[
  \hat\tau_{\mathrm{AIPW}}
  = \frac{1}{n}\sum_{i=1}^n
    \left[
      \hat\mu_1(X_i) - \hat\mu_0(X_i)
      + \frac{T_i}{\hat\pi_i}\{Y_i - \hat\mu_1(X_i)\}
      - \frac{1-T_i}{1-\hat\pi_i}\{Y_i - \hat\mu_0(X_i)\}
    \right].
\]

\medskip\noindent
\textbf{Fixed-offset augmented-model estimator.}  Within each arm,
regress $Y_i - \hat\mu_t(X_i)$ on $c_{ti}$ without an intercept to
obtain $\hat\gamma_t$, set
$\hat\mu_t^{\mathrm{aug}}(X_i) = \hat\mu_t(X_i) + \hat\gamma_t c_{ti}$,
and compute
$\hat\tau_{\mathrm{aug}}
 = n^{-1}\sum_i\{\hat\mu_1^{\mathrm{aug}}(X_i)
   - \hat\mu_0^{\mathrm{aug}}(X_i)\}$.

\medskip\noindent
\textbf{Improved augmented-model estimator.}  Within each arm,
regress $Y_i$ on $(1, \hat\mu_t^{(0)}(X_i), c_{ti})$ by OLS to
obtain $(\hat\alpha_t, \hat\beta_t, \hat\gamma_t)$, set
$\hat\mu_t^{\mathrm{aug}^+}(X_i)
 = \hat\alpha_t + \hat\beta_t\hat\mu_t(X_i) + \hat\gamma_t c_{ti}$,
and compute
$\hat\tau_{\mathrm{aug}^+}
 = n^{-1}\sum_i\{\hat\mu_1^{\mathrm{aug}^+}(X_i)
   - \hat\mu_0^{\mathrm{aug}^+}(X_i)\}$.

\subsection{Results}
\label{w10:subsec:lab-results}

Repeat the simulation $B = 2000$ times and record bias, variance,
and MSE for each estimator.  Table~\ref{w10:tab:sim} reports the
results; the improved estimator $\hat\tau_{\mathrm{aug}^+}$ is
expected to match or outperform $\hat\tau_{\mathrm{aug}}$ in both
scenarios.

\begin{table}[h]
\centering
\caption{Simulation results ($n=1000$, $B=2000$ replications,
\texttt{set.seed(2025)}).  PS model correctly specified in both
scenarios.  Entries are bias, variance, and MSE $\times 10^{3}$.
$\hat\tau_{\mathrm{aug}}$: fixed-offset augmented model.
$\hat\tau_{\mathrm{aug}^+}$: improved augmented model
\eqref{w10:eq:aug-model}.}
\label{w10:tab:sim}
\renewcommand{\arraystretch}{1.3}
\begin{tabular}{llrrrr}
\hline
Scenario & & $\hat\tau_{\mathrm{HT}}$ &
             $\hat\tau_{\mathrm{AIPW}}$ &
             $\hat\tau_{\mathrm{aug}}$ &
             $\hat\tau_{\mathrm{aug}^+}$ \\
\hline
\multirow{3}{*}{1: OR correct}
  & Bias$\times 10^3$     & $-0.37$ & $0.19$          & $0.19$          & $0.20$ \\
  & Var$\times 10^3$      & $7.78$  & $4.25$          & $4.25$          & $4.24$ \\
  & MSE$\times 10^3$      & $7.78$  & $4.25$          & $4.25$          & $4.24$ \\[4pt]
\multirow{3}{*}{2: OR misspecified}
  & Bias$\times 10^3$     & $2.41$  & $2.20$          & $2.20$          & $1.06$ \\
  & Var$\times 10^3$      & $7.72$  & $7.70$          & $7.54$          & $4.49$ \\
  & MSE$\times 10^3$      & $7.73$  & $7.71$          & $7.55$          & $4.49$ \\
\hline
\end{tabular}
\end{table}

\noindent The results reveal a striking pattern.

\begin{itemize}
  \item \textbf{Scenario~1 (OR correct).}  All four estimators are
    approximately unbiased.  $\hat\tau_{\mathrm{HT}}$ has roughly
    twice the MSE of the other three because it discards the outcome
    model entirely.  $\hat\tau_{\mathrm{AIPW}}$, $\hat\tau_{\mathrm{aug}}$,
    and $\hat\tau_{\mathrm{aug}^+}$ are essentially identical
    (MSE $\approx 4.25 \times 10^{-3}$), confirming that when the OR
    model is correctly specified the augmented parameters
    $(\hat\alpha_t, \hat\beta_t, \hat\gamma_t)$ converge to
    $(0, 1, 0)$ and the improved model incurs no efficiency cost.

  \item \textbf{Scenario~2 (OR misspecified).}  All four estimators
    remain nearly unbiased (PS is correct throughout), confirming
    double robustness.  The efficiency ordering is
    \[
      \hat\tau_{\mathrm{aug}^+} \;\ll\;
      \hat\tau_{\mathrm{aug}} \;<\;
      \hat\tau_{\mathrm{AIPW}} \;\ll\;
      \hat\tau_{\mathrm{HT}}.
    \]
    The simulation suggests that when the propensity score model is
    correct but the initial outcome model is misspecified, the
    augmented-model estimator can be more efficient than the plug-in
    AIPW estimator.  The reason is not that AIPW still carries
    first-order bias---all three PS-adjusted estimators are
    approximately unbiased---but that the targeted regression fit
    can reduce residual variance while preserving the same
    bias-correction property enforced by the IBC condition.
    The fixed-offset estimator $\hat\tau_{\mathrm{aug}}$ achieves
    only a modest 2\% MSE reduction relative to AIPW
    ($7.55$ vs $7.71 \times 10^{-3}$): with $\hat\beta_t$ fixed
    at~1, the regression can only reduce the component of the
    omitted quadratic $0.5X_i^2$ that is captured by $\hat\pi_i^{-1}$
    alone.  The improved estimator $\hat\tau_{\mathrm{aug}^+}$,
    by contrast, achieves a \emph{42\% MSE reduction} relative to AIPW
    ($4.49$ vs $7.71 \times 10^{-3}$).  The free coefficient
    $\hat\beta_t$ recalibrates the scale of the initial fit, so the
    regression of $Y_i$ on $(\hat\mu_t^{(0)}(X_i), \hat\pi_i^{-1})$
    approximates a much larger portion of $\E(Y\mid T=t, X)$, leaving
    residuals far smaller than those from the misspecified
    $\hat\mu_t^{(0)}$ alone.  The result is a substantially lower
    residual variance and hence lower MSE.
\end{itemize}

% ============================================================
\section{The Efficient Influence Function and Semiparametric Efficiency}
\label{w10:sec:eif}
% ============================================================

\emph{This section states the semiparametric efficiency result and
interprets it in light of the estimator development above; no new
derivation is given.  A complete proof requires explicit
characterization of the nuisance tangent space and is beyond the
scope of these notes; see \citet{tsiatis2006semiparametric} for a
rigorous development.}

\medskip

We now return to the efficient influence function introduced in
Chapter~9 and show that it is exactly the estimating function
underlying AIPW.  This identification is crucial: it explains why
AIPW is not only doubly robust, but also semiparametrically efficient
when both nuisance functions are estimated consistently at suitable
rates.

Under the nonparametric model for $O=(X,T,Y)$, the efficient
influence function for the ATE is
\begin{align}
  \varphi_{\mathrm{eff}}(O)
  = \frac{T}{\pi(X)}\{Y-\mu_1(X)\}
    - \frac{1-T}{1-\pi(X)}\{Y-\mu_0(X)\}
    + \mu_1(X) - \mu_0(X) - \tau.
  \label{w10:eq:eif-ate}
\end{align}
Comparing \eqref{w10:eq:eif-ate} with \eqref{w10:eq:dr-score}, the
AIPW estimating function \emph{is} the efficient influence function.
Consequently, AIPW is not merely doubly robust: it is the canonical
regular asymptotically linear estimator for the ATE when the nuisance
functions are estimated consistently at suitable rates.  The
augmentation space $\Lambda$ in \eqref{w10:eq:aug-space} corresponds
to the nuisance tangent space in semiparametric theory, and projection
onto $\Lambda^\perp$ corresponds to removing nuisance variation; see
Chapter~9, Remark~\ref{w9:rem:eif-ate} for a complementary
discussion.

\begin{thm}{Semiparametric Efficiency Bound for the ATE}
\label{w10:thm:bound}
Among regular asymptotically linear estimators of the ATE under the
nonparametric model, every such estimator has asymptotic variance at
least $\E\{\varphi_{\mathrm{eff}}(O)^2\}$.  An estimator whose
influence function equals $\varphi_{\mathrm{eff}}$ achieves this
lower bound.
\end{thm}

\begin{remark}[EIF as the semiparametric score]
\label{w10:rem:eif-as-score}
The efficient influence function plays the same role in semiparametric
theory that the score plays in parametric models: its variance gives
the information lower bound, analogous to the Cramér--Rao bound.
The three derivations in this chapter---bias correction
(Section~\ref{w10:sec:prediction}), optimal augmentation
(Section~\ref{w10:sec:class}), and the present semiparametric
identification---converge on the same estimating function, which is
the deepest explanation of why AIPW occupies a central place in the
causal inference toolkit.
\end{remark}

% ============================================================
\section{Asymptotic Inference with Estimated Nuisance Functions}
\label{w10:sec:asymp}
% ============================================================

A fundamental question in the practical use of the AIPW estimator
is: under what conditions is the effect of estimating the nuisance
functions $\pi$ and $\mu_t$ asymptotically negligible, so that
$\hat\tau_{\mathrm{AIPW}}$ behaves as if the nuisance functions were
known?  Answering this question determines when plug-in nuisance
estimators, including flexible machine-learning methods, can be used
without invalidating large-sample inference on $\tau$.

\subsection{Asymptotic normality}
\label{w10:subsec:asymp-normality}

Under suitable regularity conditions,
\[
  \sqrt{n}(\hat\tau_{\mathrm{AIPW}}-\tau)
  = \frac{1}{\sqrt{n}}\sum_{i=1}^n\varphi_{\mathrm{eff}}(O_i)
    + o_p(1)
  \overset{d}{\longrightarrow}
  N\!\bigl(0,\;\E\{\varphi_{\mathrm{eff}}(O)^2\}\bigr).
\]
The $o_p(1)$ remainder captures the effect of nuisance estimation
error.  Neyman orthogonality implies that this error enters the AIPW
estimator only through a second-order remainder: the Gateaux
derivative of the estimating equation with respect to each nuisance
function, evaluated at the truth, is zero.  A sufficient condition
for the remainder to be negligible is the \emph{product-rate
condition}
\begin{align}
  \|\hat\pi-\pi\|\cdot\|\hat\mu_t-\mu_t\| = o_p(n^{-1/2}),
  \qquad t=0,1,
  \label{w10:eq:product-rate}
\end{align}
where the norms are mean-squared-error norms.  This allows each
nuisance estimator, in principle, to converge at rate $n^{-1/4}$,
which is much weaker than the parametric rate $n^{-1/2}$ required of
each individually.  In modern applications with flexible
machine-learning methods, this statement is usually paired with
\emph{sample splitting} or \emph{cross-fitting}; without such
devices, additional empirical-process conditions are needed to justify
the asymptotic expansion.  Chapter~11 develops the cross-fitting
approach in detail.

\begin{remark}[Neyman orthogonality]
\label{w10:rem:neyman}
The first-order insensitivity of the AIPW estimating function to
nuisance perturbations is an instance of \emph{Neyman orthogonality},
a property shared by a broad class of semiparametric estimators.
Chapter~11 exploits this property systematically through
cross-fitting, which eliminates the need for additional
empirical-process conditions on the nuisance estimators.
\end{remark}

\subsection{Variance estimation and confidence intervals}
\label{w10:subsec:variance}

Because $\hat\tau_{\mathrm{AIPW}}$ is asymptotically linear with
influence function $\varphi_{\mathrm{eff}}$, its asymptotic variance
can be estimated consistently by the empirical variance of the
plug-in influence values.  Compute
\begin{align}
  \hat\varphi_i
  = \frac{T_i}{\hat\pi(X_i)}\{Y_i-\hat\mu_1(X_i)\}
    - \frac{1-T_i}{1-\hat\pi(X_i)}\{Y_i-\hat\mu_0(X_i)\}
    + \hat\mu_1(X_i)-\hat\mu_0(X_i)
    - \hat\tau_{\mathrm{AIPW}},
  \label{w10:eq:ifhat}
\end{align}
and set
\begin{align}
  \hat V
  = \frac{1}{n(n-1)}\sum_{i=1}^n(\hat\varphi_i-\bar\varphi)^2,
  \qquad
  \bar\varphi = \frac{1}{n}\sum_{i=1}^n\hat\varphi_i.
  \label{w10:eq:var-aipw}
\end{align}
The Wald confidence interval is
$\hat\tau_{\mathrm{AIPW}} \pm z_{1-\alpha/2}\sqrt{\hat V}$.

\begin{prop}{Consistency of the Variance Estimator}
\label{w10:prop:var-consistency}
If $\sqrt{n}(\hat\tau_{\mathrm{AIPW}}-\tau)
 = n^{-1/2}\sum_i\varphi_{\mathrm{eff}}(O_i)+o_p(1)$, then
$\hat V$ in \eqref{w10:eq:var-aipw} is consistent for the asymptotic
variance of $\hat\tau_{\mathrm{AIPW}}$.
\end{prop}

\begin{warn}{Extreme Propensity Scores}
\label{w10:warn:overlap}
When $\hat\pi(X_i)$ is near $0$ or $1$, the IPW weights become large
and can destabilize the estimator.  Practical remedies include overlap
diagnostics, truncation of extreme propensity scores, or restricting
the target population to a subgroup with adequate overlap.  The
augmented model approach of Section~\ref{w10:sec:ibc} provides a
complementary strategy by building the propensity score into the
outcome regression itself.
\end{warn}

% ============================================================
\section{Comparison of Regression, IPW, and AIPW}
\label{w10:sec:comparison}
% ============================================================

\begin{center}
\renewcommand{\arraystretch}{1.3}
\begin{tabular}{lcclp{4.8cm}}
\hline
Estimator & Uses $\mu_t$ & Uses $\pi$ & Consistent if & Main weakness \\
\hline
Regression (prediction) & Yes & No  & outcome model correct
  & Sensitive to OR misspecification \\
IPW (Horvitz--Thompson) & No  & Yes & propensity model correct
  & Unstable under weak overlap \\
AIPW                    & Yes & Yes & either model correct
  & Requires estimating both nuisances and careful inference (cross-fitting or regularity conditions) \\
\hline
\end{tabular}
\end{center}

\begin{remark}[When both models are correctly specified]
\label{w10:rem:both-correct}
When both models are correct, AIPW achieves the semiparametric
efficiency bound.  When overlap is weak, large IPW weights may
destabilize the estimator; the weighted regression and augmented
model approaches of Section~\ref{w10:sec:ibc} provide more stable
alternatives by incorporating balance constraints directly into the
outcome fitting step.
\end{remark}

% ============================================================
\section{Chapter Summary}
\label{w10:sec:summary}
% ============================================================

\begin{enumerate}
  \item The prediction estimator is biased when the outcome model is
    misspecified; the bias can be estimated with the propensity score
    and subtracted to yield the AIPW estimator
    (Section~\ref{w10:sec:prediction}).
  \item The AIPW estimator has two key properties.  First, double
    robustness (Theorem~\ref{w10:thm:dr}): it is consistent whenever
    either the outcome model or the propensity score model is correctly
    specified.  Second, within the class of augmented IPW estimators
    indexed by $b_t(x)$, the choice $b_t^*(x)=\E\{Y(t)\mid X\}$
    minimizes variance when both nuisance components are correctly
    specified (Theorems~\ref{w10:thm:cond-unbiased}--\ref{w10:thm:opt-b}).
  \item The optimality of AIPW follows also from a Hilbert-space
    projection: the optimal estimator is the projection of the HT
    estimator onto $\Lambda^\perp$, and the Pythagorean identity
    governs the variance reduction (Theorem~\ref{w10:thm:projection}).
  \item Double robustness can also be enforced within the
    outcome-regression fitting step, either by IPW-weighted regression
    or by an augmented regression model that includes the inverse
    propensity score as a clever covariate.  The latter is closely
    related to the targeting idea used in TMLE
    \citep{vanderLaan2006targeted} (Section~\ref{w10:sec:ibc}).
  \item The simulation study (Section~\ref{w10:sec:lab}) shows that
    when the propensity score model is correct, IPW, AIPW, and the
    augmented-model estimator are all consistent.  When the outcome
    regression is misspecified, the augmented-model estimator can
    outperform the plug-in AIPW estimator because the targeted
    regression fit may reduce residual variance while preserving the
    calibration equations.
  \item The AIPW estimating function is the efficient influence
    function for the ATE, and its variance gives the semiparametric
    efficiency bound (Section~\ref{w10:sec:eif}).
  \item For asymptotic inference, the central issue is whether the
    effect of estimating nuisance functions is asymptotically
    negligible.  Neyman orthogonality yields a second-order remainder,
    and the product-rate condition \eqref{w10:eq:product-rate} is
    sufficient for root-$n$ inference.  In modern machine-learning
    settings, cross-fitting is often used to make these conditions
    more plausible; Chapter~11 develops this approach
    (Section~\ref{w10:sec:asymp}).
\end{enumerate}

% ------------------------------------------------------------
%  Key objects
% ------------------------------------------------------------
\begin{tcolorbox}[defstyle, title={Key Objects in This Chapter}]
\begin{center}
\renewcommand{\arraystretch}{1.4}
\begin{tabular}{lp{8.5cm}}
\hline
\textbf{Symbol} & \textbf{Meaning} \\
\hline
$\tau$               & ATE $= \E\{Y(1)-Y(0)\}$ \\
$\mu_t(x)$           & Outcome regression $\E(Y\mid T=t,X=x)$ \\
$\pi(x)$             & Propensity score $P(T=1\mid X=x)$ \\
$b_t(x)$             & Control function; optimal choice $b_t^*(x)=\E\{Y(t)\mid X=x\}$ \\
$\hat\tau_{\mathrm{AIPW}}$ & AIPW estimator \eqref{w10:eq:aipw} \\
$\Lambda$            & Augmentation space \eqref{w10:eq:aug-space} \\
$\Lambda^\perp$      & Orthogonal complement of $\Lambda$ \\
$\varphi_{\mathrm{eff}}(O)$ & Efficient influence function \eqref{w10:eq:eif-ate} \\
$\hat\pi_i^{-1}$     & Clever covariate; its normal equation enforces the IBC condition \eqref{w10:eq:ibc-cond1} \\
$\hat\gamma_1$       & Targeting coefficient; $\hat\gamma_1\overset{p}{\to}0$ when outcome model correct \\
IBC                  & Internal bias calibration conditions \eqref{w10:eq:ibc-cond1}--\eqref{w10:eq:ibc-cond2} \citep{firth1998robust} \\
\hline
\end{tabular}
\end{center}
\end{tcolorbox}

% ============================================================
\section*{Exercises}
\label{w10:sec:exercises}
% ============================================================

\begin{enumerate}

\item \textbf{Bias of the prediction estimator.}
  \begin{enumerate}
    \item Verify the bias formula \eqref{w10:eq:pred-bias} by
      writing $\tau_{\mathrm{pred}}(m)-\tau$ as a difference of
      working-model errors and simplifying.
    \item Under what condition on the propensity score model does
      $\widehat{\mathrm{Bias}}(\hat\tau_{\mathrm{pred}})$ have zero
      expectation?  Provide a careful argument using iterated
      expectations.
    \item Construct a simple example (binary $X$, binary $T$,
      binary $Y$) in which both the estimated bias is nonzero and
      the outcome model is misspecified, but the AIPW estimator
      is still consistent.
  \end{enumerate}

\item \textbf{The AIPW class and optimal control functions.}
  \begin{enumerate}
    \item Verify that $\hat\tau_b$ with $b_t=\mu_t$ equals the
      AIPW estimator \eqref{w10:eq:aipw}.
    \item Using \eqref{w10:eq:var-b}, show that
      $b_t^*(x)=\E\{Y(t)\mid x\}$ minimizes conditional variance
      by completing the square in $b_t$.
    \item Suppose $\pi(x)=1/2$ for all $x$.  Simplify
      \eqref{w10:eq:var-b} and interpret the result.
  \end{enumerate}

\item \textbf{Double robustness.}
  \begin{enumerate}
    \item Verify Case~1 of Theorem~\ref{w10:thm:dr}: with
      $\mu_t=\mu_t^*$ and arbitrary $\pi$, show
      $\E\{\phi\}=0$.
    \item Verify Case~2: with $\pi=\pi^*$ and arbitrary $\mu_t$,
      show $\E\{\phi\}=0$.
    \item Provide a counterexample showing $\E\{\phi\}\neq 0$ when
      both models are misspecified.
  \end{enumerate}

\item \textbf{Projection and the Pythagorean identity.}
  \begin{enumerate}
    \item Verify $\mathrm{Cov}(\hat\theta_{\mathrm{opt}},\hat b^*)=0$
      from the definition of $\hat b^*$.
    \item Deduce \eqref{w10:eq:pythagoras} from the decomposition
      $\hat\theta_0 = \hat\theta_{\mathrm{opt}} + \hat b^*$.
    \item Explain why $\hat\theta_{\mathrm{opt}}$ is still unbiased
      for $\theta$ even though $\hat b^* \in \Lambda$ is subtracted.
  \end{enumerate}

\item \textbf{Augmented model and the clever covariate.}
  \begin{enumerate}
    \item Let $\hat\mu_1^{(0)}(X_i)$ be any initial outcome model
      fit.  Show that the normal equation for $\hat\gamma_1$ in
      the augmented model \eqref{w10:eq:aug-model} is exactly the
      IBC condition \eqref{w10:eq:ibc-cond1}, and conclude
      that \eqref{w10:eq:aug-pred} can always be written in AIPW
      form as in \eqref{w10:eq:aug-aipw}.
    \item Explain why $\hat\gamma_1 \overset{p}{\to} 0$ when the
      initial model $\hat\mu_1^{(0)}$ is correctly specified, and
      what this implies about the efficiency cost of including the
      clever covariate.
    \item Explain in words why $\hat\pi_i^{-1}$ is called the clever
      covariate: what bias does it absorb, and why does this make the
      prediction estimator a debiased estimator even when
      $\hat\mu_1^{(0)}$ is misspecified?
  \end{enumerate}

\item \textbf{Semiparametric efficiency.}
  \begin{enumerate}
    \item Show $\E\{\varphi_{\mathrm{eff}}(O)^2\}
      \leq \E\{\varphi_{\mathrm{IPW}}(O)^2\}$ by writing
      $\varphi_{\mathrm{IPW}} = \varphi_{\mathrm{eff}}
       + (\varphi_{\mathrm{IPW}}-\varphi_{\mathrm{eff}})$ and
      showing the cross term vanishes.
    \item Interpret the efficiency gain
      $\E\{\varphi_{\mathrm{IPW}}^2\}-\E\{\varphi_{\mathrm{eff}}^2\}$
      in terms of the variance of the outcome regression.
    \item Give one reason why an efficient estimator based on
      $\varphi_{\mathrm{eff}}$ may still be disfavored in
      practice relative to a simpler, less efficient alternative.
  \end{enumerate}

\end{enumerate}

% Bibliography (standalone only)
\ifSubfilesClassLoaded{%
  \bibliographystyle{plainnat}
  \bibliography{causal_inference}
}{}

\end{document}